
\SGAInlineRefTarget{sga3:xvii:section:6:site:da2ef89f51}\SGASection{6}{Extensions of a Unipotent Group by a Group of
Multiplicative Type}
\SGARefTarget{sga3:xvii:section:6}

\SGASubsection{6.1}{Statement of the Theorem}
\SGARefTarget{sga3:xvii:subsection:6:1}

\SGAStatement{\SGAInlineRefTarget{sga3:xvii:theorem:6-1-1}Theorem 6.1.1}{Let \(k\) be a field, \(U\) a
unipotent algebraic \(k\)-group, \(H\) a \(k\)-group of
multiplicative type, and \(E\) an algebraic \(k\)-group which is
an extension of \(U\) by \(H\), so that there is an exact sequence
\[
  1\longrightarrow H\longrightarrow E\longrightarrow U
    \longrightarrow1.
\]
Then the extension \(E\) is trivial and there exists a unique
lift of \(U\) in \(E\) in each of the following cases:
\begin{enumerate}
\item[(A)] \SGARefTarget{sga3:xvii:theorem:6:1:1:item:A}
  The group \(U\) is smooth and one of the following conditions
  holds:
  \begin{enumerate}
  \item[(i)] \SGARefTarget{sga3:xvii:theorem:6:1:1:item:A:i}
    \(U\) is connected and the canonical morphism
    \(E\longrightarrow U\) admits a section.
  \item[(ii)] \SGARefTarget{sga3:xvii:theorem:6:1:1:item:A:ii}
    \(U\) has a composition series whose successive quotients
    are isomorphic to \(\mathbb G_a\).
  \item[(iii)] \SGARefTarget{sga3:xvii:theorem:6:1:1:item:A:iii}
    \(H\) is étale.
  \item[(iv)] \SGARefTarget{sga3:xvii:theorem:6:1:1:item:A:iv}
    \(k\) is perfect.
  \end{enumerate}
\item[(B)] \SGARefTarget{sga3:xvii:theorem:6:1:1:item:B}
  \(U=\boldsymbol\alpha_p\) and \(k\) is perfect.
\item[(C)] \SGARefTarget{sga3:xvii:theorem:6:1:1:item:C}
  \(E\) is commutative and \(k\) is perfect.
\end{enumerate}}
\SGARefTarget{sga3:xvii:theorem:6:1:1}

\SGASubsection{6.2}{Proof of \SGARefLink{sga3:xvii:theorem:6:1:1:item:A}{6.1.1(A)}}
\SGARefTarget{sga3:xvii:subsection:6:2}

We first establish three lemmas.

\SGAStatement{Lemma 6.2.1}{Let \(S\) be a prescheme, \(E\) an
\(S\)-group prescheme which is an extension of an \(S\)-group
prescheme \(U\), with connected fibres, by an \(S\)-group \(H\)
of multiplicative type and finite type (that is, \(U\) is the
quotient of \(E\) by \(H\) for the fpqc topology).  Then \(E\) is
a central extension.}
\SGARefTarget{sga3:xvii:lemma:6:2:1}

\SGAProof{Proof}
Since \(H\) is commutative, \(U\) acts on \(H\) by inner
automorphisms, through an \(S\)-morphism of groups
\[
  u:U\longrightarrow\operatorname{Aut}_{S\text{-}\mathrm{gr}}(H).
\]
The functor
\(\operatorname{Aut}_{S\text{-}\mathrm{gr}}(H)\) is representable
by an étale \(S\)-scheme (\SGARefLink{sga3:x:corollary:5-10}{Exposé X, 5.10}), and consequently its
identity section is both an open and a closed immersion.  Since
\(U\) has connected fibres, it follows that \(u\) is the unit
morphism.

\SGAStatement{Lemma 6.2.2}{With the notation of \SGARefLink{sga3:xvii:theorem:6-1-1}{6.1.1}, if \(E\)
is trivial, there exists a unique lift of \(U\) in \(E\).}
\SGARefTarget{sga3:xvii:lemma:6:2:2}

\SGAProof{Proof}
Let \(U'\) and \(U''\) be two lifts of \(U\) in \(E\).  To prove
that \(U'=U''\), we may suppose that \(k\) is algebraically closed,
and it is enough to prove that \(\mathrm H^1(U,H)=0\).  If \(U\)
is connected, then \(U\) centralizes \(H\) by \SGARefLink{sga3:xvii:lemma:6:2:1}{6.2.1}, and hence
\[
  \mathrm H^1(U,H)
    =\operatorname{Hom}_{k\text{-}\mathrm{gr}}(U,H)=0
\]
by 2.4(ii).

In the general case, let \(U_1'\) denote the unique lift in \(E\)
of the identity component \(U^0\) of \(U\), and put
\[
  N=\operatorname{Norm}_{E}(U_1').
\]
If \(g\in E(k)\), then
\(\operatorname{int}(g)U_1'\) is a lift of \(U^0\), and is
therefore equal to \(U_1'\); thus \(N(k)\supset E(k)\).
Moreover, \(N\) contains \(H\), by \SGARefLink{sga3:xvii:lemma:6:2:1}{6.2.1}, as well as \(U_1'\),
and it follows at once that \(N=E\).  On passing to the quotient
by \(U_1'\), we are reduced to the case in which \(U\) is étale.
Since \(k\) is algebraically closed, in this case
\(\mathrm H^1(U,H)\) identifies with the ordinary cohomology
group
\[
  \mathrm H^1\bigl(U(k),H(k)\bigr),
\]
and is therefore zero, since \(U(k)\) is a finite \(p\)-group and
\(H(k)\) is uniquely \(p\)-divisible.

\SGAStatement{Corollary 6.2.3}{To prove that \(E\) is trivial, it
is enough to prove that it becomes trivial after a finite
separable extension of the base field.  In particular, we may
suppose that \(H\) is diagonalizable.}
\SGARefTarget{sga3:xvii:corollary:6:2:3}

\SGAStatement{Lemma 6.2.4}{In order that every central extension
\(E\) of \(U\) by a diagonalizable \(k\)-group \(H\) be trivial,
it is enough that every central extension of \(U\) by
\(\mathbb G_m\) be trivial.}
\SGARefTarget{sga3:xvii:lemma:6:2:4}

\SGAProof{Proof}
Induction on \(r\) first shows that the hypothesis implies that
\(E\) is trivial when \(H=\mathbb G_m^r\).  In the general case,
\(H\) embeds in \(\mathbb G_m^r\) for a suitable integer \(r\)
(this is immediate by duality); put
\[
  H''=\mathbb G_m^r/H.
\]
We obtain the exact sequence (Appendix I, \SGARefLink{sga3:i:subsection:2-1}{2.1})
\[
  \mathrm Z^1(U,H'')
    \longrightarrow\operatorname{Ext}_{\mathrm{alg}}(U,H)
    \longrightarrow
      \operatorname{Ext}_{\mathrm{alg}}(U,\mathbb G_m^r)=0,
\]
where \(U\) acts trivially on \(H\), \(H''\), and
\(\mathbb G_m^r\).  But
\[
  \mathrm Z^1(U,H'')
    =\operatorname{Hom}_{k\text{-}\mathrm{gr}}(U,H'')=0
\]
by 2.4(ii), and hence
\(\operatorname{Ext}_{\mathrm{alg}}(U,H)=0\).

\SGAProof{Proof of \SGARefLink{sga3:xvii:theorem:6:1:1:item:A}{6.1.1(A)}\SGARefLink{sga3:xvii:theorem:6:1:1:item:A:i}{(i)}}
Since \(E\longrightarrow U\) admits a section, \(E\) defines an
element \(f\) of \(\mathrm H^2(U,H)\) (Appendix I, \SGARefLink{sga3:i:definition:3-1}{3.1}).  We must
prove that \(f\) is zero.  For this it is enough to prove that a
\(2\)-cocycle
\[
  f:U\times_kU\longrightarrow H
\]
is a constant morphism, which follows from the next lemma.

\SGAStatement{Lemma 6.2.5}{Let \(U\) be a smooth connected
unipotent algebraic \(k\)-group, and \(H\) a \(k\)-group of
multiplicative type.  Then every \(k\)-morphism of preschemes
\[
  f:U\longrightarrow H
\]
is constant.}
\SGARefTarget{sga3:xvii:lemma:6:2:5}

\SGAProof{Proof}
We may suppose that \(k\) is algebraically closed.  We argue by
induction on \(\dim U\).  If \(\dim U>0\), then \(U\) has a
composition series (see \SGARefLink{sga3:xvii:lemma:3-9}{3.9})
\[
  1\longrightarrow U'\longrightarrow U
    \xrightarrow{\ \pi\ }U''\longrightarrow1,
  \qquad
  U'\simeq\mathbb G_a,\quad \dim U''<\dim U.
\]
It is enough to show that \(f\) factors through \(U''\).  Since
the graph of the equivalence relation defined by \(\pi\) is
smooth over \(k\), and hence reduced, it is enough to show that
if \(x,y\in U(k)\) have the same image \(z\) in \(U''(k)\), then
\(f(x)=f(y)\).  Now \(\pi^{-1}(z)\) is isomorphic to
\(\mathbb G_a\), so the restriction of \(f\) to
\(\pi^{-1}(z)\) factors through a reduced irreducible component
of \(H\), and hence through a \(k\)-scheme isomorphic to
\(\mathbb G_m^r\).  It remains only to note that every morphism
from \(\mathbb G_a\) to \(\mathbb G_m^r\) is constant, since
every invertible regular function on \(\mathbb G_a\) is
constant.

\SGAProof{Proof of \SGARefLink{sga3:xvii:theorem:6:1:1:item:A}{6.1.1(A)}\SGARefLink{sga3:xvii:theorem:6:1:1:item:A:ii}{(ii)}}
By \SGARefLink{sga3:xvii:lemma:6:2:1}{6.2.1}, \SGARefLink{sga3:xvii:corollary:6:2:3}{6.2.3}, and \SGARefLink{sga3:xvii:lemma:6:2:4}{6.2.4}, we may suppose that
\(H=\mathbb G_m\).  By hypothesis, \(U\) has a composition series
\[
  U\supset U_1\supset\cdots
\]
such that \(U_i/U_{i+1}\) is isomorphic to \(\mathbb G_a\).  Let
\(E_1\) be the inverse image of \(U_1\) in \(E\).  By induction
on \(\dim U\), we may suppose that the extension \(E_1\) is
trivial; let \(U_1'\) be the unique lift of \(U_1\).  Arguing as
in the normalizer argument of \SGARefLink{sga3:xvii:lemma:6:2:2}{6.2.2}, we find that \(U_1'\) is
invariant in \(E\).  After taking the quotient by \(U_1'\), we
are reduced to the case \(U=\mathbb G_a\).  The \(S\)-scheme
\(E\), where \(S=U\), is then a torsor under the \(S\)-group
\(\mathbb G_m\times_kS\), and therefore admits a section, since
\(\operatorname{Pic}(S)=0\) (\SGARefLink{sga3:viii:corollary:4:3}{Exposé VIII, 4.3}).  The extension
\(E\) is then trivial by (A)\SGARefLink{sga3:xvii:theorem:6:1:1:item:A:i}{(i)}.

\SGAProof{Proof of \SGARefLink{sga3:xvii:theorem:6:1:1:item:A}{6.1.1(A)}\SGARefLink{sga3:xvii:theorem:6:1:1:item:A:iii}{(iii)} (\(U\) smooth, \(H\) étale)}
Suppose first that \(U\) is connected.  The group
\(U_{\overline{k}}\) then has a composition series whose
successive quotients are isomorphic to \(\mathbb G_a\), so
\(E_{\overline{k}}\) is trivial by (A)\SGARefLink{sga3:xvii:theorem:6:1:1:item:A:ii}{(ii)}.  Since \(H\) is
étale, the unique lift of \(U_{\overline{k}}\) in
\(E_{\overline{k}}\) is plainly the identity component of
\(E_{\overline{k}}\); this lift is therefore already defined
over \(k\).  In the general case, after quotienting by the
identity component of \(E\), we are reduced to the case where
\(E\) is étale, and then to the case where \(E\) is completely
split (\SGARefLink{sga3:xvii:corollary:6:2:3}{6.2.3}).  Since \(U(k)\) is a \(p\)-group and \(H(k)\) has
order prime to \(p\), we may take the Sylow \(p\)-subgroup of
\(E(k)\) as a lift of \(U(k)\).

\SGAProof{Proof of \SGARefLink{sga3:xvii:theorem:6:1:1:item:A}{6.1.1(A)}\SGARefLink{sga3:xvii:theorem:6:1:1:item:A:iv}{(iv)} (\(U\) smooth, \(k\) perfect)}
If \(U\) is connected, it has a composition series whose
successive quotients are isomorphic to \(\mathbb G_a\)
(\SGARefLink{sga3:xvii:lemma:4:1:2:item:b}{4.1.2(b)}), and we apply (A)\SGARefLink{sga3:xvii:theorem:6:1:1:item:A:ii}{(ii)}.  In the general case, the
preceding argument reduces us to the case in which \(U\) is
étale, and then to the case in which \(U\) is completely split
and \(H\) is diagonalizable (\SGARefLink{sga3:xvii:corollary:6:2:3}{6.2.3}).  Using a characteristic
composition series of \(H\),
\[
  H\supset H^0\supset{}_{F^n}(H)\supset0,
\]
we are reduced to the case in which \(H\) is of one of the
following three types:
\begin{enumerate}
\item[(a)] \(H\) is étale;
\item[(b)] \(H=\mathbb G_m^r\);
\item[(c)] \(H\) is radicial.
\end{enumerate}
In case \SGARefLink{sga3:xvii:appendix:A}{(a)} we apply (A)\SGARefLink{sga3:xvii:theorem:6:1:1:item:A:iii}{(iii)}, and in case \SGARefLink{sga3:xvii:appendix:C}{(c)} we apply \SGARefLink{sga3:xvii:lemma:1:6}{1.6}.
Finally, in case \SGARefLink{sga3:xvii:section:b}{(b)}, Hilbert's \SGARefLink{sga3:viii:corollary:4:5}{Theorem 90} shows that
\(E\longrightarrow U\) admits a section, so it is enough to show
that
\[
  \mathrm H^2(U,\mathbb G_m^r)
    =\mathrm H^2\bigl(U(k),\mathbb G_m^r(k)\bigr)=0.
\]
But \(U(k)\) is a finite \(p\)-group, whereas
\(\mathbb G_m^r(k)\) is uniquely \(p\)-divisible, since \(k\) is
perfect.

\SGASubsection{6.3}{Proof of \SGARefLink{sga3:xvii:theorem:6:1:1:item:B}{6.1.1(B)} and (C)}
\SGARefTarget{sga3:xvii:subsection:6:3}

By \SGARefLink{sga3:xvii:lemma:6:2:1}{6.2.1}, \SGARefLink{sga3:xvii:corollary:6:2:3}{6.2.3}, and \SGARefLink{sga3:xvii:lemma:6:2:4}{6.2.4}, it is enough to prove (B) when
\(H=\mathbb G_m\).  We therefore have an exact sequence
\[
  1\longrightarrow\mathbb G_m\longrightarrow E
    \longrightarrow\boldsymbol\alpha_p\longrightarrow1.
\]
Since \(\mathbb G_m\) is smooth, it gives an exact sequence of
restricted Lie algebras (Appendix II, \SGARefLink{sga3:xvii:proposition:3:2}{3.2})
\begin{equation}
  0\longrightarrow\operatorname{Lie}\mathbb G_m
    \longrightarrow\operatorname{Lie}E
    \longrightarrow\operatorname{Lie}\boldsymbol\alpha_p
    \longrightarrow0.
  \tag{*}\label{sga3:xvii:subsection:6-3:equation:star}
\end{equation}
The group \(\boldsymbol\alpha_p\) has height \(1\), so the
extension \(E\) is trivial if and only if the exact sequence
\((*)\) of restricted Lie algebras splits (Appendix II, \SGARefLink{sga3:ii:subsection:2-2}{2.2}).
The algebra \(\operatorname{Lie}\boldsymbol\alpha_p\) is
generated by an element \(X\ne0\) such that \(X^{(p)}=0\)
(Appendix II, \SGARefLink{sga3:xvii:proposition:2:1}{2.1}).  It is therefore enough to show that \(X\)
lifts to an element \(Z\) of \(\operatorname{Lie}E\) for which
\(Z^{(p)}=0\).  By \SGARefLink{sga3:xvii:lemma:5-8-2:item:ii}{5.8.2(ii)}, there is a unipotent element \(Z\)
of \(\operatorname{Lie}E\) which lifts \(X\).  Since the unipotent
part of \(\operatorname{Lie}E\) plainly has dimension at most
\(1\), it follows necessarily that \(Z^{(p)}=0\).
